\section{Integer boundary-path diagrams and duon collision protocol}
\label{app:integer-boundary-path-diagrams}

This appendix records the diagram grammar used by the manual for retained boundary paths. It is a manual presentation of the exact row objects defined in the main note; the row entries remain integer, rational-pair, cycle, residue, sign, and topology data.

\subsection{Boundary-link row}
\label{appD:boundary-link-row}

A retained boundary link on level \(\ell\) and window \(\omega\) is recorded as
\[
D_e^{(\ell,\omega)}=
\left(
 c_e,k_e,N_e,
 \chi_e,
 Q_e,S_e,
 R_e,R_{0,e},
 T_e,T_{0,e},
 B_e^*,\widehat\Lambda_e,
 P_e,P_{0,e},
 \Sigma_e
\right).
\]
The entries record cycle data, chirality, effective-charge counts, loss counts, torsion pair, support-count witness, attenuation exponent, AOD momentum proxy, and topology signature.

\subsection{Topology signature and incidence}
\label{appD:topology-incidence}

The topology signature is
\[
\Sigma_e=(\tau_e,\sigma_e,\mathbf a_e,\partial_e,\mathcal I_e),
\]
where \(\tau_e\) is a cell or face type index, \(\sigma_e\in\{-1,0,+1\}\) is an orientation sign, \(\mathbf a_e\in\mathbb Z^m\) is a fractal or scale address, \(\partial_e\) is a boundary incidence list, and \(\mathcal I_e\) is an integer incidence matrix or tensor. Incidence signs are read as
\[
\epsilon_{ej}=\epsilon(\Sigma_e,\Sigma_j)\in\{-1,0,+1\}.
\]

\subsection{Transport flux and windowed current-flux}
\label{appD:transport-current-flux}

With \(r_e\) the centered phase representative, define
\[
\Delta\Phi_e=(r_e,N_e),
\qquad
q_e=(Q_e,S_e),
\qquad
G_{T,e}=(T_{0,e}+T_e,T_{0,e}).
\]
Let \(A_{B,e}=(A_{B,e}^{+},A_{B,e}^{-})\) and \(A_{R,e}=(A_{R,e}^{+},A_{R,e}^{-})\). The exact retained transport flux is
\[
u_e^{(\mathrm{tr})}
=
\operatorname{Rat}
\left[
 r_e Q_e A_{B,e}^{+} A_{R,e}^{+}(T_{0,e}+T_e),
 N_e S_e A_{B,e}^{-} A_{R,e}^{-}T_{0,e}
\right].
\]
The longitudinal accumulated transport is
\[
U_K^{(\ell)}(\omega)=
\sum_{e\in E_K(\omega)}\epsilon_e^{\parallel}u_e^{(\mathrm{tr})}.
\]
For \(|\omega|_{\mathrm{bip}_{\ell}}\in\mathbb Z_{>0}\), the current-flux rate is
\[
\mathcal J_K^{(\ell)}(\omega)=
\operatorname{Rat}
\left[
U_{K,\mathrm{num}}^{(\ell)}(\omega),
U_{K,\mathrm{den}}^{(\ell)}(\omega)|\omega|_{\mathrm{bip}_{\ell}}
\right].
\]

\subsection{Carrier-family partition}
\label{appD:carrier-partition}

The first retained transport examples use the disjoint partition
\[
E_K(\omega)=E_{K,\mathrm{duon}}(\omega)\cup E_{K,\mathrm{tetron}}(\omega)\cup E_{K,\mathrm{other}}(\omega),
\]
\[
E_{K,a}\cap E_{K,b}=\varnothing\qquad(a\neq b).
\]
Then \(U_K\) and \(\mathcal J_K\) split exactly into duon, tetron, and other carrier contributions.

\subsection{Transverse witness}
\label{appD:transverse-witness}

The exact transverse witness is
\[
w_{\perp,j}^{\mathrm{exact}}=
\sum_e \epsilon_{ej}^{\perp}\chi_eu_e^{(\mathrm{tr})}.
\]
For the two-sided wire witness,
\[
\epsilon_{e,+}^{\perp}=+1,
\qquad
\epsilon_{e,-}^{\perp}=-1,
\]
so
\[
w_{\perp}^{+}=\sum_e\chi_eu_e^{(\mathrm{tr})},
\qquad
w_{\perp}^{-}=-\sum_e\chi_eu_e^{(\mathrm{tr})}.
\]
The notation guard is
\[
\boxed{B^*\neq \mathcal J_K\neq w_\perp}.
\]

\subsection{Retained-boundary diagram value}
\label{appD:retained-boundary-diagram-value}

A retained-boundary or retained-path diagram value is
\[
\mathcal A_{\mathbb Z}(\mathcal G)=
(W_{\mathcal G},K_{\mathcal G},N_{\mathcal G},\chi_{\mathcal G},\Sigma_{\mathcal G}),
\]
where \(W_{\mathcal G}\) is a rational weight pair, \(K_{\mathcal G}\) is an accumulated phase residue, \(N_{\mathcal G}\) is the accumulated phase modulus, \(\chi_{\mathcal G}\) is an accumulated chirality sign, and \(\Sigma_{\mathcal G}\) is an accumulated topology signature. For a path \(\gamma\),
\[
K_\gamma=\sum_{e\in\gamma}k_e\pmod{N_\gamma}.
\]
The retained-path accumulation is
\[
\mathcal Z_{\partial\mathcal O}^{(\ell,\omega)}=
\bigoplus_{\gamma\in\Gamma_{\partial\mathcal O}(\omega)}
(W_\gamma,K_\gamma,N_\gamma,\chi_\gamma,\Sigma_\gamma).
\]

\subsection{Integer collision vertex}
\label{appD:integer-collision-vertex}

For incoming boundary signs \(a,b\in\{-1,0,+1\}\), the balanced ternary update is \(c=a\oplus_3b\). The retained collision vertex is
\[
\mathcal K_{\mathrm{coll}}^{\mathbb Z}:
D_a\times D_b\times\Sigma_{ab}
\longrightarrow
D'\oplus\Pi_{\mathrm{export}}\oplus\mathcal R_{\mathrm{cut}}.
\]
The update predicates are
\[
a\oplus_3 b,
\qquad
Q_a+Q_b,
\qquad
K_a+K_b\pmod N,
\qquad
\chi_a\chi_b,
\qquad
\epsilon(\Sigma_a,\Sigma_b).
\]
The outcome class is one of balanced, push, pull, reclosure, export, or rupture/refinement.
